This paper studies how discriminatory structure can arise among interacting
learning agents. Two points bear stating.

The model is not a claim about people. Its agents are single-layer perceptrons,
its groups are arbitrary inert labels, and its "tolerance" is a single scalar. We
use the vocabulary of in-group and out-group because it names the structure
compactly, not because the model licenses inferences about human prejudice, whose
causes are historical, institutional and material in ways nothing here
represents. Where we discuss the affective-versus-ideological polarisation
literature, we do so to point out that the model identifies a parameter that
selects between two regimes, not to adjudicate an empirical question about any
actual polity.

The language-agent study instructs agents to treat one group's statements as less
credible than another's, and measures the result. This is the intervention under
study, so it cannot be avoided, but it is confined: the groups are arbitrary
partitions with no reference to any real social category, protected or otherwise,
and the labels carry no content connected to the agenda. No human subjects are
involved and no personal data is used. We report the instruction wording in full
(Appendix~\ref{app:llm-prompts}) because reproducibility requires it and because
readers should be able to judge how strong the intervention is.

The practical direction of the work is defensive. The result implies that a
multi-agent system can occupy a discriminatory collective state while every agent
in it passes an individual audit, and that the transition into that state is
sharp in the strength of a per-agent bias. That argues for population-level
measurement of deployed multi-agent systems, and the order parameters used here
are inexpensive enough to serve as one.
